\section*{Резюме}
\addcontentsline{toc}{section}{Резюме}
\label{sec:rezyume}

Курсовата работа изследва доказателствената сила на електронен документ, подписан с
квалифициран електронен подпис (КЕП), и правилата за оспорването му в българския
граждански процес. Тя започва с правното действие на КЕП по Регламент (ЕС) №~910/2014
(eIDAS) и Закона за електронния документ и електронните удостоверителни услуги и с
механизма на подписа, като разграничава какво доказва валидирането и какво не доказва.
След това разглежда представянето на електронния документ и формалната му доказателствена
сила по ГПК, предмета и сроковете на оспорването, разпределението на доказателствената
тежест по чл.~193, ал.~3 ГПК и ограниченията по чл.~18а ЗЕДЕУУ. Последната част анализира
злоупотребите с КЕП, задължението на доставчика да провери самоличността и
представителната власт, включително по казус от практиката на Комисията за регулиране на
съобщенията, и отговорността на доставчика.

Изводът е, че при валиден КЕП тежестта да докаже неистинността пада върху лицето,
посочено като подписал, макар валидирането да не доказва кой е дал командата за
подписване. Затова решаващи са данните, които доставчикът е длъжен да пази, и презумпцията
за неговата вина при неправилно издадено удостоверение.

\medskip
\noindent\textbf{Ключови думи:} квалифициран електронен подпис; електронен документ;
оспорване на истинността; доказателствена тежест; удостоверителни услуги; eIDAS.
